\hypertarget{appendix-b60686ba06a1}{}\label{referencestart:appendix-b60686ba06a1}
\section{Fourier analysis and spectral phase}\label{appendix-b60686ba06a1-fourier-analysis-and-spectral-phase}

\hypertarget{appendix-b60686ba06a1-fourier-analysis-and-spectral-phase}{}

\textbf{Curriculum depth:} Developed · \textbf{First seen:} XR01 · \textbf{Latest development:} University courses

\begin{allmightroledefinition}

\subsection{Definition and mechanism}\label{appendix-b60686ba06a1-definition-and-mechanism}

\hypertarget{appendix-b60686ba06a1-definition-and-mechanism}{}

Fourier coordinates encode frequency magnitude and phase; phase places structure. DFT analyzes a finite periodic block, FFT computes that transform efficiently, and finite-window leakage differs from bin spacing. Interior one-sided amplitude bins need conjugate-pair normalization.

\end{allmightroledefinition}

\subsection{Prerequisites}\label{appendix-b60686ba06a1-prerequisites}

\hypertarget{appendix-b60686ba06a1-prerequisites}{}

\begin{itemize}
\tightlist
\item
  \href{../07-Vision-XR-and-Playable-Worlds-P02/main.pdf\#appendix-1077ea8fc737}{Continuous signals and discrete samples (Vision, XR, and Playable Worlds, Part 2, Reference notes: Continuous signals and discrete samples)} --- Filtering/frequency analysis operates on a declared domain and its samples.
\end{itemize}

\subsection{Used by}\label{appendix-b60686ba06a1-used-by}

\hypertarget{appendix-b60686ba06a1-used-by}{}

\begin{itemize}
\item
  \hyperref[appendix-6a7869dd79a6]{Sampling aliasing and reconstruction}
\item
  \href{../03-Geometry-Cameras-and-Image-Formation-P01/main.pdf\#appendix-d5b3c692d9f0}{Convolution and filter footprints (Geometry, Cameras, and Image Formation, Part 1, Reference notes: Convolution and filter footprints)}
\item
  \href{../07-Vision-XR-and-Playable-Worlds-P02/main.pdf\#appendix-8b76b0efd27f}{PSF OTF and MTF (Vision, XR, and Playable Worlds, Part 2, Reference notes: PSF OTF and MTF)}
\item
  \href{../07-Vision-XR-and-Playable-Worlds-P02/main.pdf\#appendix-b2c1846508bf}{Perceptual metrics and display calibration (Vision, XR, and Playable Worlds, Part 2, Reference notes: Perceptual metrics and display calibration)}
\item
  \href{../03-Geometry-Cameras-and-Image-Formation-P01/main.pdf\#appendix-3c73f7862157}{Exceptional projections and dimension (Geometry, Cameras, and Image Formation, Part 1, Reference notes: Exceptional projections and dimension)} \#\# Recall and next depth
\end{itemize}

Recover the two tested sinusoid amplitudes and explain leakage. Pending: FFT implementations, windows and finite-boundary estimates. Written coverage is not a claim of learner mastery or accepted numbered-curriculum progress.

\subsection{Source and navigation}\label{appendix-b60686ba06a1-source-and-navigation}

\hypertarget{appendix-b60686ba06a1-source-and-navigation}{}

\href{../07-Vision-XR-and-Playable-Worlds-P02/main.pdf\#04-visual-computing-ii-ch042-fourier-coordinates-reveal-what-a-filter-removes}{Expanded teaching explanation (Vision, XR, and Playable Worlds, Part 2, Ch. 9)} · \hyperref[appendix-fd090d982a15-xr01---independent-seeing-and-representation-source]{Independent coverage} · \href{../07-Vision-XR-and-Playable-Worlds-P01/main.pdf\#appendix-4ad47bdd269a}{Dependency Map (Vision, XR, and Playable Worlds, Part 1, Reference notes: Dependency Map)}

\subsection{University course development}\label{appendix-b60686ba06a1-university-course-development}

\hypertarget{appendix-b60686ba06a1-university-course-development}{}

Extend the existing frequency-domain treatment with visual interpretation and filtering consequences. A Fourier coefficient records how strongly a sinusoidal basis function occurs at a given frequency and phase. In images, the origin of a two-dimensional spectrum represents slow spatial variation; moving farther from it corresponds to finer variation. A broad bright center indicates low-frequency structure, while energy far from the center often represents edges and texture. Filtering multiplies each frequency by the filter response: a low-pass response suppresses fine detail, and a high-pass response emphasizes it. Aliasing folds unresolved high frequencies into lower ones because sampled sinusoids can share identical sample values. This explains moiré patterns and animated shimmer, and shows why a spectrum is useful for diagnosing sampling problems rather than only for theoretical signal processing.

\href{../03-Geometry-Cameras-and-Image-Formation-P01/main.pdf\#03-visual-computing-i-ch003-sampling-continuous-signals-and-image-formation}{Sampling continuous signals and image formation (Geometry, Cameras, and Image Formation, Part 1, Ch. 10)} · \href{../03-Geometry-Cameras-and-Image-Formation-P01/main.pdf\#03-visual-computing-i-ch004-image-sampling-filtering-and-efficient-reconstruction}{Image sampling, filtering, and efficient reconstruction (Geometry, Cameras, and Image Formation, Part 1, Ch. 11)} · \href{../05-Spatial-Worlds-Engine-State-and-Simulation-P01/main.pdf\#04-visual-computing-ii-ch030-edges-contours-and-subpixel-localization}{Edges, contours, and subpixel localization (Spatial Worlds, Engine State, and Simulation, Part 1, Ch. 11)}
